
Or \(\underline{N}'\) est de
rang \(d-2\), donc \(\lambda^{d-1}(N')=0\), d'où le résultat.

\subsection*{4.2.}

Soit \(i:Y\to X\) une immersion régulière. Soit \(r\) un entier \(\ge1\), et \(s:X\longrightarrow \underline{\underline{P}}_X^r\) la section définie par l'homomorphisme : \((\underline{O}_X)^{r+1}\to\underline{O}_X\), passage au quotient par le sous-faisceau engendré par les \(r\) dernières coordonnées. Si \(\underline{\underline{P}}_X^r\) est défini par la \(\underline{O}_X\)-Algèbre graduée \(\underline{O}_X[T_0,\ldots,T_r]\), l'idéal \(\underline{I}\) définissant la section \(s\) est l'idéal engendré par les sections \(T_1,\ldots,T_r\), donc \(\underline{I}/\underline{I}^2\) est isomorphe à \((\underline{O}_X[T_0])^r\), donc l'idéal correspondant de \(\underline{O}_p\) à \((\underline{O}_X)^r\).

Comme la section \(s\) est évidemment une immersion régulière, il en est de même pour \(s\circ i\), et on peut écrire la suite exacte des faisceaux conormaux correspondante (1.7). On obtient ainsi l'égalité suivante dans \(K^{\bullet}(Y)\), en notant \(\underline{N}\) le faisceau conormal de \(i\), et \(\underline{N}_1\) le faisceau conormal de \(s\circ i\) :
\[
  N_1=N+r.
\]
On voit donc que pour toute immersion régulière \(i\), on peut, en la composant avec une section : \(X\to\underline{\underline{P}}_X^2\), ou bien : \(X\to\underline{\underline{P}}_{\underline{\underline{P}}_X^1}^{1}\), obtenir une immersion régulière satisfaisant la propriété du lemme 4.1.

\textbf{Théorème 4.3.} \uline{Soient} \(i:Y\to X\) \uline{une immersion fermée régulière}, \(N\) \uline{la classe dans} \(K^{\bullet}(Y)\) \uline{du faisceau conormal de} \(Y\) \uline{dans} \(X\). \uline{On suppose que} \(X\) \uline{est quasi-compact et a un faisceau ample}. \uline{Alors}, \uline{avec les notations de} 3.3, \uline{on a pour tout} \(y\in K^{\bullet}(Y)\), \uline{et tout} \(p\ge1\) :
\begin{equation}\tag{4.3.1}
 i_*\bigl(\lambda^p(N,y)\bigr)\equiv \lambda^p(i_*(y))
 \qquad\text{\uline{mod}. \uline{torsion}},
\end{equation}
\(\lambda^p(N,y)\) \uline{étant l'élément défini dans} (V 5.3).

i) On suppose d'abord que l'hypothèse de 4.1 est satisfaite ; on va dans ce cas montrer que \(i_*\bigl(\lambda^p(N,y)\bigr)=\lambda^p(i_*(y))\). Comme \(f_*\circ f^*=\Id\), cette
